%==============================================================================
\chapter{Référentiel métier et bonnes pratiques apicoles}
%==============================================================================
Ce chapitre précise les connaissances métier qui fondent les règles de gestion,
les seuils paramétrables et les indicateurs du système. Il constitue la base
fonctionnelle sur laquelle reposent les tableaux de bord et les alertes.

\section{Facteurs de production du miel}
La production d'une colonie dépend directement de la météo, de la disponibilité
du pollen et du flux de nectar. Lorsque le flux de nectar est élevé, la reine
pond davantage, la population culmine et le rendement atteint son maximum. En
période de pénurie alimentaire, la ponte et la population diminuent. Le système
doit donc corréler la production mesurée aux périodes et aux saisons, ce qui
justifie le filtrage par mois, semaine, saison et année des tableaux de bord.

\section{Emplacement et configuration d'un site}
Le choix de l'emplacement d'un site conditionne la productivité. Les valeurs de
référence suivantes peuvent alimenter les champs et les contrôles de saisie du
module de gestion des sites.
\begin{longtable}{>{\bfseries}p{5.5cm} p{8.5cm}}
\toprule
\textbf{Critère} & \textbf{Valeur de référence} \\
\midrule
\endhead
Proximité des ressources & Sources de nectar et de pollen à environ 1 km. \\
Accès à l'eau & Eau potable accessible, plusieurs litres consommés par jour et par colonie. \\
Éloignement des habitations & 100 m en zone forestière, 200 m en zone arbustive, 300 m en zone ouverte. \\
Éloignement des cultures traitées & Au moins 3 à 4 km des cultures conventionnelles utilisant des pesticides. \\
Surélévation et inclinaison & Ruches surélevées à environ 1,5 m du sol et légèrement inclinées pour évacuer l'humidité. \\
Protection & Abri contre le soleil intense, le vent et les inondations. \\
\bottomrule
\end{longtable}

\begin{figure}[htbp]
\centering
\includegraphics[width=0.48\linewidth]{hive-003}
\caption{Exemple de ruche installée sur un site : surélévation sur pieds et planche d'envol.}
\label{fig:ruche-installee}
\end{figure}

\section{Types de ruches et rendement de référence}
\begin{longtable}{>{\bfseries}p{5.5cm} p{8.5cm}}
\toprule
\textbf{Type de ruche} & \textbf{Caractéristiques et rendement} \\
\midrule
\endhead
Ruche traditionnelle à rayons fixes & 6 à 10 kg de miel en rayons par saison, non recommandée en apiculture biologique. \\
Ruche à barrettes supérieures & Largeur standard de 32 à 35 cm, inspection indépendante des rayons. \\
Ruche à cadres (Langstroth, Est-Africaine) & Rayons à miel réutilisables plusieurs saisons après récolte. \\
\bottomrule
\end{longtable}
Ces ordres de grandeur servent de base au paramétrage des seuils de production
et à l'estimation du retour sur investissement par ruche.

\section{Classification des modèles de ruches}
On distingue deux grandes familles de ruches, dont la modélisation doit rester
ouverte pour couvrir les usages locaux et internationaux.
\begin{longtable}{>{\bfseries}p{4.5cm} p{9.5cm}}
\toprule
\textbf{Famille} & \textbf{Modèles représentatifs} \\
\midrule
\endhead
Ruches à cadres (verticales) & Dadant, Langstroth, Voirnot, Layens, Aumônière, Nationale (GB), WBC, Tonelli. \\
Ruches traditionnelles sans cadres & Warré, Kényane TBH, Tanzanienne TBH, tronc, paille. \\
Autres modèles régionaux & Lucitana, Oksman, Smith, CDB, Duvauchelle, Jarry, Congres, Zanders, modèles slovènes. \\
\bottomrule
\end{longtable}
En France, les modèles les plus courants sont les ruches Dadant, Langstroth et
Voirnot. Toutes les ruches à cadres partagent une composition commune, avec des
dimensions précises au millimètre près, ce qui conforte le modèle de données
retenu.
\begin{encadre}
\textbf{Ancrage du modèle :} le modèle de référence du projet est la ruche
Dadant, dont le corps accueille jusqu'à 10 cadres. Cette valeur fonde la
capacité maximale de 10 cadres par corps ou par hausse retenue dans les règles
de gestion. Le système reste toutefois paramétrable pour supporter d'autres
modèles.
\end{encadre}

\subsection{Composition commune d'une ruche à cadres}
\begin{longtable}{>{\bfseries}p{4cm} p{10cm}}
\toprule
\textbf{Élément} & \textbf{Rôle} \\
\midrule
\endhead
Plancher / planche d'envol & Base de la ruche et zone d'atterrissage des abeilles. \\
Corps (socle) & Compartiment principal du couvain, contient les cadres de base. \\
Hausse & Étage d'extension destiné à la récolte du miel. \\
Cadre & Support amovible des rayons de cire, installé sur un slot identifié. \\
Grille à reine & Sépare le corps des hausses et cantonne la reine dans le corps. \\
Couvre-cadre et toit & Ferment et protègent la ruche des intempéries. \\
Nourrisseur & Permet le ravitaillement en nourriture de la colonie. \\
\bottomrule
\end{longtable}

\begin{figure}[htbp]
\centering
\includegraphics[width=0.85\linewidth]{hive-001}
\caption{Coupe d'une ruche à cadres et rôle de chaque élément (toit, nourrisseur, couvre-cadre, hausse, grille à reine, corps, plancher).}
\label{fig:coupe-ruche}
\end{figure}

\section{Protocole d'inspection des colonies}
Les inspections sont réalisées de préférence les jours ensoleillés et portent
sur des points de contrôle précis qui structurent le rapport de visite.
\begin{itemize}
  \item développement du couvain (oeufs, larves, cellules operculées)
  \item remplissage des alvéoles en miel et en pollen
  \item présence de parasites ou de maladies
  \item récolte de nectar, de pollen et de propolis
\end{itemize}
\begin{encadre}
\textbf{Contrainte opérationnelle :} une visite ne doit pas dépasser 45 minutes.
Au-delà, l'agitation de la colonie augmente et se propage aux ruches voisines.
Le champ \emph{durée} du rapport de visite peut être contrôlé par rapport à
cette limite de référence.
\end{encadre}

\section{Essaimage et division des colonies}
Un essaimage imminent se signale par la présence de cellules royales sur les
bords des rayons. Quelques jours avant l'émergence d'une nouvelle reine, la
vieille reine quitte la colonie avec la moitié des ouvrières. Le système doit
donc permettre d'anticiper la division. Trois méthodes de référence sont
prévues comme motifs d'intervention.
\begin{longtable}{>{\bfseries}p{4.5cm} p{9.5cm}}
\toprule
\textbf{Méthode} & \textbf{Principe} \\
\midrule
\endhead
Formation de nucléus & Retirer 1 à 2 rayons de couvain avec abeilles vers une ruche neuve. \\
Nucléus de reine & Transférer la reine avec un rayon sans cellules royales. \\
Essaimage artificiel & Déplacer la ruche mère et placer une nouvelle ruche à l'ancien emplacement avec 2 à 3 rayons de jeune couvain. \\
\bottomrule
\end{longtable}

\section{Indicateurs de bien-être et causes de baisse}
Une colonie saine présente un nid à couvain bien développé aux différents
stades, des alvéoles remplies de ressources, une activité de butinage visible
et l'absence de parasites ou de maladies. Ces critères alimentent l'évaluation
qualitative de l'état de santé et le niveau de productivité noté de 1 à 3 dans
le rapport de visite.

La fuite ou l'abandon d'une ruche résulte de perturbations excessives ou de
conditions inadéquates (manque de nourriture ou d'eau, températures extrêmes).
Laisser une réserve de miel à la colonie lors de la récolte réduit ce risque.
Ces phénomènes justifient les alertes de chute soudaine d'effectif ou de
production et les seuils paramétrables associés.

\section{Incidence sur les seuils paramétrables}
\begin{longtable}{>{\bfseries}p{4.5cm} p{9.5cm}}
\toprule
\textbf{Paramètre système} & \textbf{Ancrage métier} \\
\midrule
\endhead
Seuil de collecte de miel & Rendement de référence par type de ruche et par saison. \\
Seuil d'alerte de baisse & Chute anormale liée à une fuite, une maladie ou une pénurie alimentaire. \\
Durée maximale de visite & Limite de 45 minutes par inspection. \\
Distances de sécurité du site & 100, 200 ou 300 m selon la zone, 3 à 4 km des cultures traitées. \\
Périodes d'analyse & Corrélation production et saison (flux de nectar). \\
\bottomrule
\end{longtable}

